% =========================
% Explanation: Experimental Design (informal, readable)
% =========================
\subsection*{Explanation of Experimental Design (informal)}
This section explains how we run the benchmarking experiments so results are statistically meaningful and comparisons are fair.
It specifies budgets, repeats, seeds, stopping rules, parallelism, and any early-stopping / multi-fidelity mechanisms.

\begin{itemize}
    \item \textbf{Evaluation unit (what counts as one trial).}
    We define precisely what ``one evaluation'' means (e.g., a full training run for a fixed number of epochs,
    a fixed token budget, or a tabular lookup).
    This matters because the primary notion of runtime in optimization benchmarking is the number of objective evaluations.
    All methods are charged identically: each evaluated configuration counts as one evaluation,
    including invalid or failed trials (handled consistently).

    \item \textbf{Budget tiers (and why multiple budgets).}
    We test at least three evaluation budgets \(B\in\{B_1,B_2,B_3\}\) (e.g., 25/50/100).
    Multiple budgets matter because optimizers can behave differently in the low-budget regime (cold start) versus higher budgets
    (refinement and exploitation). Reporting across tiers prevents cherry-picking and makes scaling behavior explicit.

    \item \textbf{Stopping rules.}
    Each optimizer run stops when it reaches its evaluation budget \(B\); no method may exceed \(B\) evaluations.
    If additional constraints are used (e.g., wall-clock caps), they are applied uniformly across methods and reported explicitly.
    If a run terminates early due to unrecoverable failures, the failure is logged and handled consistently in the analysis.

    \item \textbf{Random seeds and independent repetitions.}
    Because HPO/NAS methods are stochastic, a single run is not representative.
    We therefore run \(R\) independent repetitions per (method, benchmark instance, budget), each with a different optimizer seed.
    Performance is aggregated across these \(R\) runs (e.g., median and IQR), and statistical tests account for run-to-run variability.

    \item \textbf{Controlling training stochasticity (evaluation seeds).}
    Training itself can be stochastic (random initialization, data order, dropout, nondeterministic GPU operations).
    To improve comparability, we optionally use \emph{common random numbers} (CRN):
    the evaluation seed for a configuration \(\lambda\) is deterministically derived from the benchmark id and \(\lambda\).
    Thus, if different optimizers evaluate the same \(\lambda\), they observe the same training realization, reducing noise in comparisons.
    Optionally, final selected configurations are re-evaluated with multiple fresh seeds for more stable reporting;
    these re-evaluations are reported separately and do not count toward the HPO budget.

    \item \textbf{Initialization policy.}
    Some methods (e.g., Bayesian optimization) require an initial design before model-based proposals are meaningful.
    To keep comparisons fair, we use a shared initialization policy (e.g., Sobol or Latin hypercube for continuous dimensions
    and uniform sampling for categorical variables), allocate the same number of initial evaluations \(n_0\) when required,
    and count these evaluations toward the total budget \(B\).

    \item \textbf{Multi-fidelity / early stopping (if allowed).}
    If the benchmark supports multi-fidelity evaluation (e.g., fewer epochs, smaller dataset fraction, lower resolution),
    we specify a shared fidelity ladder and a shared scheduler policy (e.g., ASHA/Hyperband).
    In this setting, all methods follow the same resource allocation logic and differ only in how they propose configurations.
    We report both the number of evaluated configurations and the total consumed resources (e.g., total epochs/tokens),
    but define the primary runtime axis for benchmarking metrics as completed evaluations for consistency with fixed-budget/fixed-target definitions.

    \item \textbf{Parallel evaluation policy (synchronous vs.\ asynchronous).}
    Optimizers can gain unfair advantages if they exploit asynchrony or different degrees of parallelism.
    We therefore fix a parallelism policy:
    either (i) synchronous batches with a fixed worker count \(W\), or (ii) asynchronous execution with shared dispatcher semantics
    applied to all methods. The policy, worker count, and how partial results are handled are stated explicitly.
    Regardless of wall-clock behavior, our primary runtime axis for benchmarking is completed evaluation count.

    \item \textbf{Hardware/software and determinism constraints.}
    We report the framework, hardware, and determinism controls (fixed RNG seeds, deterministic kernels where supported,
    fixed data loader order). This is essential for reproducibility and for interpreting variability, especially in GPU-based training
    where strict determinism may not always be achievable.
\end{itemize}
